\chapter{Vortex-oriented \texorpdfstring{$S_4$}{S4} three-cycle lock}

Rotation audit оставил восемь orientation-preserving affine candidates:
ровно восемь three-cycles стандартного \(S_4\)-triplet. Прежний continuous
Wilson vacuum не совпадал с ними ни по оси, ни по углу. Здесь проверяется
другая возможность: может ли unit vortex winding сам ориентировать
дискретный family holonomy без выбора одной permutation вручную.

\section{Кубический инвариант стандартного triplet}

Реализуем стандартное представление как
\begin{equation}
 F=\left\{x\in\mathbb R^4:\sum_{a=1}^4x_a=0\right\},
 \qquad \|x\|=1.
\end{equation}
На нём существует канонический \(S_4\)-инвариант
\begin{equation}
 I_3(x)=\sum_{a=1}^4x_a^3.
\end{equation}
Метод множителей Лагранжа даёт точную границу
\begin{equation}
 -\frac1{\sqrt3}\leq I_3(x)\leq\frac1{\sqrt3}.
\end{equation}
Равенство достигается только на восьми ориентированных тетраэдрических
вершинах
\begin{equation}
 \pm\frac1{\sqrt{12}}(3,-1,-1,-1)
\end{equation}
и их четырёх перестановках. После отождествления \(x=Bn\), где
\(B:\mathbb R^3\to F\) --- ортонормированное вложение, это ровно четыре
неориентированные оси three-cycle rotations.

\section{Угол из характера, а не из прежнего Wilson functional}

Для \(R_n(\theta)\in SO(3)\) характер стандартного triplet равен
\begin{equation}
 \chi_3(\theta)=\Tr R_n(\theta)=1+2\cos\theta.
\end{equation}
В интервале \(0<\theta<\pi\) условие \(\chi_3=0\) имеет единственное
решение
\begin{equation}
 \theta=\frac{2\pi}{3}.
\end{equation}
Это характер three-cycle в стандартном представлении. Следовательно,
нулевой характер фиксирует нужный угол без обращения к наблюдаемым данным.

Этот шаг не является малой поправкой к прежнему continuous Wilson vacuum
\(\theta_\star\). Он определяет отдельную discrete family--defect branch.
Поэтому старые коэффициенты \(V_{\rm gap}(\theta_\star)\) нельзя переносить
в новую ветвь автоматически.

\section{Совместный неотрицательный functional}

Пусть \(r=|\Phi|\) --- нормированный модуль charge-two pairing field, а
\(\nu=\operatorname{wind}\Phi=\pm1\) --- unit vortex sector. Рассмотрим
\begin{equation}
 \begin{split}
 V_\nu(r,\theta,n)={}&(r^2-1)^2
 +r^2\bigl(1+2\cos\theta\bigr)^2\\
 &+r^2\left(1-\sqrt3\,\nu I_3(Bn)\right),
 \qquad \|n\|=1.
 \end{split}
 \label{eq:family-defect-three-cycle-lock}
\end{equation}
Каждое слагаемое неотрицательно. Поэтому глобальный минимум равен нулю и
достигается тогда и только тогда, когда
\begin{equation}
 r=1,
 \qquad
 \theta=\frac{2\pi}{3},
 \qquad
 I_3(Bn)=\frac{\nu}{\sqrt3}.
\end{equation}

При \(\nu=+1\) остаются четыре oriented axes, соответствующие четырём
three-cycles одной ориентации. При \(\nu=-1\) получаются четыре обратных
цикла. Объединение двух unit-winding sectors даёт в точности все восемь
кандидатов, оставленных предыдущим \(SO(3)\)-filter.

\begin{theorem}[Vortex-oriented three-cycle selection]
Functional~\eqref{eq:family-defect-three-cycle-lock} имеет ровно четыре
глобальных минимума в каждом фиксированном unit-winding sector. Их family
holonomies являются четырьмя \(S_4\) three-cycles ориентации \(\nu\).
Следовательно, знак топологического winding заменяет ручной выбор между
cycle и inverse cycle.
\end{theorem}

\section{Hessian и Majorana kernel}

В нормированных координатах Hessian в каждом минимуме имеет eigenvalues
\begin{equation}
 \lambda_r=8,
 \qquad
 \lambda_\theta=6,
 \qquad
 \lambda_{n,1}=\lambda_{n,2}=6.
\end{equation}
Особенно важно, что угловая и обе осевые stiffness совпадают без настройки
relative coefficient. Для каждого минимума
\begin{equation}
 K_n=\frac{2\pi}{3}[n]_\times
\end{equation}
имеет rank два и одномерное ядро. Поэтому прежний exact-one mechanism
сохраняется:
\begin{equation}
 \dim\ker K_n=1,
 \qquad
 3\longrightarrow1.
\end{equation}
Максимальные numerical residuals полного exhaustive scan меньше
\(8\times10^{-16}\).

\section{Сверка с литературной архитектурой}

Superconnection formalism действительно объединяет gauge connection и
пространственно зависимый mass/tachyon field и применяется к interface и
Callias-type index problems; соответствующая современная сверка дана в
\path{arXiv:2106.01591}. Локализованные winding-mass zero modes обсуждаются
также в \path{arXiv:1905.09315}. Возможность фиксировать flavour alignment
граничными условиями реализована в orbifold constructions, например в
\path{arXiv:1910.04175}. Non-Abelian vortices с внутренними orientational
moduli дают физический образ family axis на defect, см.
\path{arXiv:1612.09419}.

Эти работы подтверждают допустимость архитектурных ингредиентов, но не
выводят specific functional~\eqref{eq:family-defect-three-cycle-lock}.
Положительный результат настоящего gate является project-internal
invariant theorem, а не заимствованной моделью.

\section{Что прошло и что осталось}

Новый functional одновременно:
\begin{itemize}
 \item выбирает ненулевой normalized condensate \(r=1\);
 \item фиксирует exact three-cycle angle \(2\pi/3\);
 \item выбирает тетраэдрическую family axis;
 \item связывает orientation cycle с vortex winding;
 \item сохраняет один Majorana core mode.
\end{itemize}

Но пока не доказано, что весь functional возникает как supertrace одной
локальной boundary-superconnection curvature. В частности, единичный
radial scale, constant term в \((r^2-1)^2\) и знак coupling
\(\nu I_3\) должны следовать из parent grading, topological term или
spectral normalization, а не задаваться после результата.

\begin{nogocriterion}[Supertrace-origin gate]
Нельзя считать three-cycle lock полным parent-action выводом, пока
functional~\eqref{eq:family-defect-three-cycle-lock} не получен из одного
заранее фиксированного graded curvature или boundary spectral action.
Добавление трёх его слагаемых как независимых phenomenological potentials
возвращает прежнюю проблему ручных sector weights.
\end{nogocriterion}

Таким образом, discrete selector и exact-one defect mechanism впервые
имеют общий coefficient-free zero locus. Следующий узкий расчёт --- найти
минимальную graded superconnection, квадрат которой порождает radial,
character и cubic-orientation terms с общей нормировкой.

Следующий projector-supercurvature gate уточняет эту запись. После замены
\(y=\nu x\) четыре axes выбираются независимо от winding, а \(\nu\)
ориентирует угол \(\pm2\pi/3\). Radial и axis conditions затем
факторизуются в одну shifted-square equation; поэтому текущая сумма трёх
terms является zero-locus representative, а не минимальной формой action.

Машинный exhaustive ledger сохранён в
\path{s2t_v4_family_defect_three_cycle_lock_gate_results.json}; генератор
проверки ---
\path{s2t_v4_family_defect_three_cycle_lock_gate.py}.